%% Abstract variants, one per target venue.
%%
%% main.tex carries the PRIMARY version below.  To retarget the manuscript,
%% replace the body of its abstract environment with one of the others and
%% re-run scripts/../abswc against the stated limit.
%%
%% This file is reference material and is not \input by main.tex.

%% ------------------------------------------------------------------------
%% PRIMARY -- Journal of the Royal Society Interface
%% limit 200 words, this version 198
%% ------------------------------------------------------------------------
\noindent
Evolutionary game theory usually identifies the payoff an individual earns with
the quantity that governs how often it is copied. Ecosystems of deployed
artificial agents break that identification: an agent competes in a development
race, but whether its design spreads is decided by a principal who does not play
it and who bears only a fraction of the harm the design causes. We call the
difference between the two payoff functionals the selection wedge, embed it in a
four-strategy model of an artificial intelligence race, and analyse the
replicator and finite-population dynamics. The wedge collapses to one control
parameter, an effective liability, in which every invasion condition is affine,
so all thresholds are closed form. Three consequences follow. Screening designs
before deployment removes only a weakly dominated design and leaves the long-run
unsafe frequency unchanged. That frequency is not monotone in liability: over an
intermediate window liability taxes the retaliation on which conditional safety
depends, the unsafe state becomes uninvadable, and crossing the window's lower
edge lifts unsafe behaviour from zero to about a third of encounters. Coupling
enforcement to diffusing capability produces hysteresis of exactly computable
width. Governance binds at the selection layer, not the design layer.

%% ------------------------------------------------------------------------
%% FALLBACK -- PNAS Nexus / Dynamic Games and Applications
%% limit 250 words, this version 249
%% ------------------------------------------------------------------------
\noindent
Evolutionary game theory usually identifies the payoff an individual earns with
the quantity that governs how often it is copied. Ecosystems of deployed
artificial agents break that identification: an agent competes in a development
race, but whether its design spreads is decided by a principal who does not play
it and bears only a fraction of the harm the design causes. We call the
difference between the two payoff functionals the selection wedge, embed it in a
four-strategy model of an artificial intelligence race, and analyse the
replicator and finite-population dynamics. The wedge collapses to one control
parameter, an effective liability $L$, in which every invasion condition is
affine, so all thresholds are closed form. Four consequences follow. First,
screening designs before deployment removes only a weakly dominated design and
leaves the long-run unsafe frequency unchanged. Second, unconditional and
conditional safety are exact payoff twins, so the barrier protecting safe play
is set by a neutrally drifting composition variable and falls there by a factor
of $77.7$. Third, long-run harm is not monotone in $L$: over an intermediate
window liability taxes the retaliation conditional safety depends on, the unsafe
state becomes uninvadable, and crossing its lower edge lifts the unsafe
frequency from $0$ to $0.32$. Fourth, coupling enforcement to diffusing
capability produces hysteresis of computable width. Assortative matching narrows
the window and closes it; misattributed blame widens it as $1/(1-q)$. Governance
binds at the selection layer, not the design layer, and some liability regimes
are worse than either extreme.

%% ------------------------------------------------------------------------
%% ORIGINAL -- technical-report version, no venue
%% limit none words, this version 503
%% ------------------------------------------------------------------------
\noindent
We study evolutionary dynamics in which the payoff earned in the game and the
payoff that governs replication are different functions of the same
interaction. The separation is the defining feature of an ecosystem of
deployed AI agents: an agent acts in a competitive race, but whether its
\emph{design} is replicated is decided by a principal who does not play that
race and who bears only a fraction $\lambda$ of the external harm the design
causes. We call the difference between the two functionals the selection
wedge, embed it in the reduced four-strategy AI race of Fern\'andez Domingos
and Han, and analyse the resulting replicator and finite-population dynamics.
The wedge collapses to a single control parameter, the effective liability
$L=\lambda h$, and every invasion condition is affine in $L$, which yields
closed-form thresholds. Four consequences follow. First, a conditionally
unsafe design weakly dominates the unconditionally unsafe one at every $L$, so
a pre-deployment filter that scores designs against a safe environment removes
only a weakly dominated design and, under a common start measure, leaves the
long-run unsafe frequency unchanged. Second, unconditional and conditional
safety are exact payoff twins, so the invasion barrier of the safe face is set
by a neutrally drifting composition variable, and it falls by a factor of
$77.7$ across that face. Third, long-run harm is \emph{not monotone} in
liability: on $L\in(4.27,42.77)$ liability taxes the retaliation that
conditional safety relies on, the unsafe coexistence becomes uninvadable, and
crossing the lower edge from below lifts the basin-averaged unsafe frequency
from $0$ to $0.32$. A window of this kind exists at every prize and
private-risk level we examine, with multiplicative width between $4.2$ and
$12.8$, and in a finite population the unsafe attractor survives between $65$ and
$3.9\times10^{13}$ generations. Fourth, coupling the dynamics to a
non-decreasing stock of diffused capability that erodes enforcement produces
hysteresis whose width is exactly computable and depends on the erosion
channel only through the enforcement that survives full diffusion, not on the
speed of diffusion. We then test how far the mechanisms travel: they are
inherited by any selection dynamics driven by the sign of a payoff difference,
they survive non-linear liability schedules and richer conditional designs,
and the valley closes only when execution error exceeds about $7.5\%$ -- at
which point the protected state has stopped being safe. Two features of the
environment then move the valley in opposite directions by computable amounts.
Assortative matching shrinks it entirely from above, because the lower edge is
exactly independent of assortment, and closes it at an assortment of $0.354$.
Booking a fraction $q$ of the blame for an interaction against the wrong party
widens it as $1/(1-q)$, because the design that a first strike exploits is by
construction the one with a spotless record; every other measurement error we
consider, including zero-mean noise, flat charges and systematic
under-reporting, is absorbed exactly. The results locate the
governance problem at the selection layer rather than at the design layer, and
identify liability regimes that are worse than either extreme.

